\hypertarget{process-and-privilege-model}{%
\chapter{2. Process and privilege
model}\label{process-and-privilege-model}}

Chapter 1 set out the kernel facilities and how the confinement is woven
from them. A kennel does not assemble itself: a small set of processes
builds each one and supervises it while it runs, and the reference
monitor's tamperproof property turns, in mechanism, on which of those
processes holds privilege and for how long. The shape of that set is
where the detail begins, which processes exist while a kennel runs, what
privilege each carries, and where the single privileged step sits. The
design register's \texttt{split-the-uid} reduces, here, to one question
asked of every process on the host: which file capability does it hold,
and for how long.

The answer almost everywhere is none. Of the processes that exist while
a workload runs, exactly one holds host-elevated privilege, it holds it
for the length of one construction, and then it is gone.

\hypertarget{the-processes}{%
\section{2.1 The processes}\label{the-processes}}

A running kennel is a small set of processes, and naming them is most of
the privilege story.

\texttt{kennel} is the operator's entry point, a 13-line shim
(\texttt{kennel-shim}) that holds no authority of its own: it detects
which environment it is in and execs the right libexec, the host-side
CLI (\texttt{kennel-cli}) on the host or the in-cage spawn path inside a
kennel. \texttt{kenneld} is the per-user supervisor, socket-activated by
\texttt{systemd\ -\/-user} on the first \texttt{kennel\ run} and
resident for the session, one per logged-in user. It runs as the user
with no capabilities. It validates and executes the settled policy for
each kennel, drives construction, takes binder node 0, drains the audit
ringbuf, and reaps the workload, and it does all of it unprivileged.
That a daemon this central holds no capability is the first half of
\texttt{split-the-uid} in mechanism: \texttt{kenneld} has the operator's
identity and none of root's authority, and the authority it needs for
the one privileged step it does not hold but borrows, briefly, from a
separate binary.

\texttt{kennel-privhelper} is that binary, and the only one on the host
with elevated privilege. Beneath it sit the processes that hold no host
privilege. \texttt{kennel-bin-init} is the kennel's PID 1: a trusted
root-owned binary, trapped inside the new namespaces with no host
capability, that supervises the facades and the workload, reaps the
kennel's processes, and restarts a facade if one dies. It holds PID 1 so
that no untrusted workload does, keeping that position, with its reaping
duties and its host-facing reparent, on a trusted process rather than on
the code the kennel exists to confine. The rest hold no host privilege:
the per-kennel \texttt{host-netproxy} and the SSH bastion are plain user
processes, the in-kennel facades are confined workload-side code, and
the workload itself is dropped to the operator with its bounding set
cleared. \texttt{host-netproxy} is a separate process for an
architectural reason rather than a privilege one: it is the post-policy
dialer, doing the blocking, adversarial network I/O once
\texttt{kenneld} has made the decision, so that I/O never sits inside
the supervisor (\texttt{control-not-data-plane}). The same eviction is
made on the other adversarial input a session produces: the workload's
terminal output, with its attacker-shaped ANSI escapes, is parsed by the
unprivileged \texttt{kennel} CLI and never at the daemon, which sees the
PTY as opaque bytes to route. Between them, network I/O and terminal
parsing are the two data-plane jobs kept out of the supervisor, which is
what leaves it a pure control plane. The privilege-bearing surface of
the whole system is one small binary and three smaller sub-helpers,
described next; everything else runs as the user.

\hypertarget{one-privileged-component}{%
\section{2.2 One privileged component}\label{one-privileged-component}}

\texttt{kennel-privhelper} is installed with the file capabilities
\texttt{cap\_setuid}, \texttt{cap\_setgid}, \texttt{cap\_setfcap},
\texttt{cap\_sys\_admin=ep}, setuid root at mode \texttt{4755} only as a
fallback for filesystems that cannot carry capability xattrs. Each of
the four earns its place and no more is taken. \texttt{cap\_setuid},
\texttt{cap\_setgid}, and \texttt{cap\_setfcap} cover writing the
identity map and dropping to the operator; \texttt{cap\_setfcap} is
needed because the map's \texttt{0\ 0\ 1} line and the operator line are
written in a single \texttt{write(2)}. \texttt{cap\_sys\_admin} is the
kernel's gate on writing a \texttt{uid\_map} that maps host uid 0, and
it covers the mount and \texttt{pivot\_root} construction of the view.
The factory holds nothing else: no \texttt{cap\_net\_admin}, no BPF
capability.

The rarer privileged acts are not folded into the factory's capability
set; they are delegated to three single-purpose sub-helpers the factory
execs only when a policy needs them, each carrying one capability for
one job. \texttt{kennel-privhelper-net} holds \texttt{cap\_net\_admin}
and mirrors the kennel's loopback prefixes onto the host's \texttt{lo}
over netlink, bringing the same \texttt{/28} and \texttt{/64} up on both
sides so the host can present a kennel's inbound listeners at the
kennel's own address; the two are identical by construction, not
connected. \texttt{kennel-privhelper-bpf} holds \texttt{cap\_bpf},
\texttt{cap\_net\_admin}, \texttt{cap\_perfmon} and, in host network
mode only, attaches the egress BPF to the kennel's cgroup.
\texttt{kennel-privhelper-mounts} holds \texttt{cap\_sys\_admin} and
over-mounts the exclusive-bind sentinel. This is \texttt{rule-of-1} and
\texttt{quarantine-the-unsafe} carried into privilege itself: the
host-elevated surface is partitioned so that the capability for a narrow
job lives in a binary that can do only that job, and holding
\texttt{cap\_net\_admin} there grants the one scoped act, not arbitrary
network administration (T2.4).

\hypertarget{the-factory-privilege-bounded-to-one-construction}{%
\section{2.3 The factory: privilege bounded to one
construction}\label{the-factory-privilege-bounded-to-one-construction}}

The privhelper is a kennel factory. Its one provisioning op,
\texttt{construct}, does all of the privileged work for a kennel in a
single short-lived invocation; the only standalone op beside it is
\texttt{del-addr}, the address removal at teardown.

\texttt{kenneld} invokes \texttt{kennel-privhelper\ construct} over a
\texttt{SOCK\_SEQPACKET} socketpair and sends the construction half of
the Plan: the uid and gid maps, the loopback configuration and the
per-kennel addresses to add, the binderfs parameters, the view bind
list, the pivot target, the egress BPF payload as a framed tail, and any
stdio or pty descriptors over \texttt{SCM\_RIGHTS}. The privhelper
parses this host-side before it holds any namespace, and parses it as
hostile: the operator line of the map is the caller's real uid taken
from \texttt{/proc} ownership and never read from the wire, and the
request is refused if it names an address outside the kennel's allocated
IPv4 \texttt{/28} or IPv6 \texttt{/64}, a cgroup the caller does not
own, a gid the caller is not in, or any uid but the caller's own. The
decoder for that half is bounded and fuzzed, because a bug in it is the
one bug at this layer that matters (T5.4).

Only after validation does privilege touch the kernel. The factory holds
neither \texttt{cap\_net\_admin} nor the BPF caps, so it drives the
host-\texttt{lo} loopback mirror and the host-mode egress BPF attach by
exec'ing the narrow sub-helpers; it provenance-checks and opens the
root-owned \texttt{kennel-bin-init} by descriptor, and then
\texttt{clone3}s a child with \texttt{CLONE\_NEWUSER} \textbar{}
\texttt{NEWNS} \textbar{} \texttt{NEWPID} \textbar{} \texttt{NEWIPC}
\textbar{} \texttt{INTO\_CGROUP}, adding \texttt{NEWNET} in every mode
but host, so the child is PID 1 of a fresh namespace set, born directly
inside the kennel's cgroup rather than migrated into it afterward. The
migration was not free: profiling added to the spawn path in 0.3.0
traced 10 to 13 ms of otherwise-idle stall to the post-hoc move into the
cgroup, and removing it by cloning into the cgroup is most of what
brings a kennel up in under 10 ms from the command line. The dynamic
spawn an agent triggers directly is faster again, an observed 3.5 to 6.5
ms depending on the hardware, because it skips the CLI and the process
startup that comes with it, and that margin is what makes an agent's
spawn-use-reap loop viable. It also closes the window in which a
freshly-migrated process runs briefly unaccounted, so the cgroup's
limits hold from the first instruction. The user namespace is
operator-owned: the child clones as the operator and self-escalates to
the kennel's uid 0 to do construction, which is precisely what lets the
unprivileged \texttt{kenneld} reach the new binderfs instance later
through \texttt{/proc/\textless{}init\textgreater{}/root}. In that child
the factory writes the maps, brings up the in-namespace loopback, builds
the view, mounts binderfs and chowns the device to the operator,
\texttt{pivot\_root}s and detaches the host root, and \texttt{fexecve}s
\texttt{kennel-bin-init} by the descriptor it opened before the pivot,
with empty argv and envp, because the host path to that binary no longer
exists inside the view.

Then the factory exits. It is not a supervisor and not a reaper proxy;
it reports the init's host pid and is gone. \texttt{kennel-bin-init},
now PID 1 of its own namespace, outlives it, reparents to
\texttt{kenneld} (which marked itself a child subreaper at startup), and
\texttt{kenneld} waits on it directly for the workload's exit status. No
privileged process is resident for the life of a kennel, and none is
shared across kennels. The privilege is bounded to the single
map-writing construction and then withdrawn, which is the conservative
half of a deliberate trade: a long-running privileged daemon owning the
same capabilities would cost fewer execs and buy continuous privileged
exposure, and the current design refuses that exposure. The factory
child is the only transient uid-0 actor in the system, and it never runs
with host root visible to it, so the worst a construction-half decoder
bug can be is a host-root bug, which is why that decoder is the most
heavily checked code at this layer (T5.4).

\hypertarget{the-trusted-init-and-the-drop}{%
\section{2.4 The trusted init and the
drop}\label{the-trusted-init-and-the-drop}}

\texttt{kennel-bin-init} is the root-owned binary the factory execs as
PID 1, and it is deliberately small (338 SLOC) and
\texttt{\#!{[}forbid(unsafe\_code){]}}. It makes no policy decision and
runs no mount, netlink, device, or filesystem-lookup code; that work was
done by the factory before the pivot. Its path comes from the root-owned
deployment configuration; the privhelper verifies the binary is
root-owned and not group- or other-writable, opens it before the
\texttt{clone3}, and \texttt{fexecve}s it by that descriptor. Exec by
descriptor is the only way in, not an optimisation: after the
\texttt{pivot\_root} the host path the binary came from is gone from the
view, so there is nothing left to exec by name, and the descriptor pins
the exact file vetted before the pivot, with no window in which a
symlink or a swapped path could redirect it. The root ownership is
tamperproofing rather than power: root confers no host capability inside
the kennel, but a root-owned binary is one the operator-uid workload
cannot rewrite, and PID 1 running as the kennel's uid 0 is one the
operator cannot signal or \texttt{ptrace}. It is launched into a cage
where its root means only that the workload cannot touch it. Once
running, it pulls the supervision half of the Plan from \texttt{kenneld}
over the binder bus, then forks and supervises the facades and the
workload.

The workload is dropped. It runs as the operator uid, never as the
kennel's uid 0; only \texttt{kennel-bin-init} holds that, trapped
post-pivot with no ambient host capability (T3.1).
\texttt{PR\_SET\_NO\_NEW\_PRIVS} is set unconditionally before the
workload runs, the bounding set is cleared per policy, Landlock is
sealed, the cgroup BPF is attached, and \texttt{setrlimit} caps are
applied. Against the classic escalation, a setuid binary executed inside
the kennel, the framework's \texttt{exec} invariants
\texttt{deny\_setuid}, \texttt{deny\_setgid}, and \texttt{deny\_setcap}
hold categorically, so a workload cannot climb back up through a binary
it runs (T3.1). There is identity here, the operator's uid, and there is
no authority attached to it that the policy did not grant: the second
half of \texttt{split-the-uid}, now at the bottom of the process tree.

\hypertarget{two-identities-not-one}{%
\section{2.5 Two identities, not one}\label{two-identities-not-one}}

That arrangement is where Kennel parts company with the common rootless
container. A rootless container maps the invoking user to uid 0 inside
its user namespace and runs the workload as root there; the confinement
is the namespace boundary, and inside it the workload is root, with
nothing between it and whatever the namespace still reaches. Kennel's
user namespace carries two identities instead. The \texttt{0\ 0\ 1} line
maps host root to the kennel's uid 0, and that root is the trusted side:
the construction child that self-escalates to it to mount and pivot, and
\texttt{kennel-bin-init} that holds it as PID 1. The operator line maps
the operator to itself, and that is the workload, unprivileged, never
uid 0.

Two identities in one namespace means the ordinary Unix DAC line between
root and a user is kept inside the kennel rather than collapsed. The
trusted components are root and the workload is not, so the process the
kennel exists to confine cannot write a root-owned path, cannot chown a
file it does not own, cannot do any of what the discretionary model
reserves to uid 0, and it is stopped by ownership alone, before
Landlock, seccomp, or the namespaces are consulted. The maps are precise
identity lines, not a subuid range, and the granted gids are mapped the
same way; but in a namespace whose only identities are root and a
non-root user, gids are ornamental, since the privilege boundary is the
uid line and nothing hangs off group membership. What stands inside the
kennel is the separation Unix has drawn between a privileged owner and
an unprivileged user for as long as it has existed, with the framework
on the privileged side and the workload on the other.

The inversion is deliberate. Rootless buys the workload root inside a
box and trusts the box; Kennel keeps a real root and a real non-root
user in the same namespace and puts the workload on the unprivileged
side of a line the kernel has enforced for decades. The namespace,
Landlock, and seccomp are the confinement above; this is the floor
beneath them, and it is the one layer that asks for no new mechanism,
because it is only uid 0 declining to be the workload's.

\hypertarget{the-unsafe-at-this-layer}{%
\section{2.6 The unsafe at this layer}\label{the-unsafe-at-this-layer}}

The privilege model is also where the volume's quarantined
\texttt{unsafe} mostly lives. Two of the five unsafe-bearing crates are
the kernel-interface surface of the privilege model:
\texttt{kennel-lib-syscall} (946 SLOC), which holds the raw syscalls,
the seccomp filter, and the namespace primitives, and
\texttt{kennel-lib-landlock} (249 SLOC), which holds the Landlock
ruleset construction; \texttt{kennel-lib-bpf} (811 SLOC), the third, is
the loader the factory and \texttt{kenneld} use to attach the egress
program. Each is small, single-purpose, in the TCB, and the only place
its kind of \texttt{unsafe} is permitted to appear
(\texttt{quarantine-the-unsafe}).

The point worth drawing out is that privilege and \texttt{unsafe} are
quarantined separately, and neither coincides with the binaries one
would expect. \texttt{kenneld} (7319 SLOC) and
\texttt{kennel-privhelper} (1764) are the two privileged anchors of the
system, and both carry no \texttt{unsafe} of their own: the privhelper's
power is file capabilities, and the few \texttt{unsafe} calls its
construction needs are borrowed from \texttt{kennel-lib-syscall}. So the
largest privileged binary holds no \texttt{unsafe}, and the crates that
hold \texttt{unsafe} hold no privilege of their own. A reader auditing
for memory hazard and a reader auditing for privilege are reading two
different, small, named sets of crates (\texttt{read-by-the-hostile}),
and that \texttt{kenneld}, the central daemon and the thing an attacker
most wants, is both unprivileged and memory-safe-by-default is the shape
the rest of the volume builds on (T5.3).

The processes and their privilege are the static picture. How they talk
to each other while a kennel runs is the next mechanism the design left
abstract, and it has a single answer on Linux: the binder.
